\chapter{Glossary}
\label{app:glossary}

This glossary defines domain-specific and thesis-specific terms used regularly throughout the chapters, in the sense they carry in this thesis specifically.

\begin{description}
  \item[Programme] A degree programme at TU Wien; this thesis covers the Bachelor Informatics and Master Software Engineering programmes specifically (Chapter~\ref{chap:methodology}, Section~\ref{sec:overall-scope}).
  \item[Exam Subject] The second level of the curriculum hierarchy (Programme $\to$ Exam Subject $\to$ Module $\to$ Course), a named grouping of modules within a programme (Section~\ref{sec:curriculum-abstraction}).
  \item[Module] The third level of the curriculum hierarchy; a named unit containing one or more courses, carrying a single obligation category (Section~\ref{sec:curriculum-abstraction}).
  \item[Course (course item)] The fourth, most granular level of the curriculum hierarchy; refers to any course-item entry regardless of teaching format (VO, UE, VU, PR, etc.).
  \item[Obligation category / obligation type] Whether a module is mandatory, constrained elective, or fully elective (Section~\ref{sec:obligation}); the underlying institutional taxonomy differs between the Bachelor and Master programmes (Section~\ref{sec:course-obligation-abstraction}).
  \item[Planning status] The lifecycle state of a course in a student's plan: not planned, parked, planned, or done (Section~\ref{sec:lifecycle}).
  \item[Parking Stage] The intermediate holding area for shortlisted courses that have not yet been assigned to a semester (Section~\ref{sec:parking}).
  \item[Prototype] The Figma prototype used to elicit requirements in the formative study (Chapter~\ref{chap:formative-study}), distinct from the implemented tool this thesis's Design and Implementation chapters describe.
  \item[Theme (C-identifier)] One of the 64 consolidated, design-relevant need categories produced by grouping the initial interview codes (Chapter~\ref{chap:formative-study}, Section~\ref{sec:theme-consolidation}); cited in running text by its identifier, e.g.\ (C21). Each theme carries the interviews and verbatim excerpts that informed it. The full set is given in Section~\ref{app:codebook}.
  \item[Code] An individual, inductively developed label attached to a coded excerpt during initial coding (Chapter~\ref{chap:formative-study}, Section~\ref{sec:initial-coding}). The 435 initial codes are the input to theme consolidation and are not cited individually in this thesis; the themes they were grouped into are.
  \item[Feature (FEATURE-XXX)] One of the 13 features derived from themes in Chapter~\ref{chap:from-themes-to-requirements}, each carrying a set of numbered requirements.
  \item[Requirement (Req.)] A single numbered requirement under a feature. Within Chapter~\ref{chap:from-themes-to-requirements}, a feature's own heading uses its name directly (e.g.\ FEATURE-005); from Design onward, a feature is cited by its cross-reference label, e.g.\ Feature~\ref{feat:005}, Req.~Y, so the citation always resolves to the feature's current number even if features are renumbered.
  \item[OOS (Out of Scope)] One of the four system-wide scope boundaries (OOS1--OOS4); OOS1 and OOS2 are general and stated in the Methodology chapter (Section~\ref{sec:overall-scope}), OOS3 and OOS4 are specific and stated in Chapter~\ref{chap:from-themes-to-requirements} (Section~\ref{sec:overall-scope-refined}) alongside the findings that justify them.
  \item[Nested Model] Munzner's nested model for visualisation design and validation~\cite{munzner2009}, used as this thesis's umbrella structuring framework (Chapter~\ref{chap:methodology}, Section~\ref{sec:paradigm}).
  \item[Table View] The primary planning canvas: semesters as vertical columns containing course cards, organised for commitment-stage planning (Section~\ref{sec:table-spatial}).
  \item[Graph View] The exploratory canvas: the curriculum rendered as a left-to-right hierarchical layout, organised for browsing and sensemaking rather than planning (Section~\ref{sec:graph-ltr}).
  \item[Recommendation Panel] The independently-toggleable left-side panel showing per-family course recommendations with inspectable evidence (Section~\ref{sec:rec-panel}).
  \item[Compliance Engine] The backend rule engine invoked after every planning action, returning fulfilment status, ECTS statistics, hard violations, and soft warnings (Section~\ref{sec:compliance}).
  \item[Dashboard] The persistent right-panel component showing compliance and progress information at four aggregation levels (Section~\ref{sec:dashboard}).
  \item[Semester Picker] The menu used to assign a course to a semester or the Parking Stage from wherever it is triggered (Section~\ref{sec:semester-picker}).
\end{description}

\TDmajor{Four glossary entries are missing}{(1)~\textbf{Dependency}, the word this thesis most needs to pin down, has no entry, and the text uses it in three senses: the containment relation the graph draws, the formal or hard prerequisite relation the curriculum publishes and the compliance engine blocks on, and the recommended or soft ordering the engine holds as an advisory. State all three and say which one a bare ``dependency'' means. (2)~\textbf{ECTS}, used 124 times and expanded nowhere. (3)~\textbf{StEOP}, used eight times and expanded nowhere. (4)~\textbf{Scenario~A / Scenario~B}, the two conditions, referred to by name throughout Chapters~\ref{chap:evaluation} and~\ref{chap:discussion}.}

\begin{description}\item[]
\end{description}
